\section{Discussion}
\label{sec:discussion}

\subsection{Key Findings and Their Implications}

Our results establish three findings with distinct implications for calibration research and LLM deployment practice.

\textbf{Finding 1: Adversarial construction method determines calibration outcome.} Human-adversarial NLI examples (AdvGLUE MNLI: $\Delta$ECE $= +0.071$) produce calibration degradation; model-in-loop adversarial NLI examples (ANLI R1/R2: $\Delta$ECE $< 0$) produce calibration improvement. This distinction --- absent in prior adversarial robustness work, which pools all ``adversarial'' conditions --- is the key methodological contribution of our experimental design. For practitioners evaluating LLM deployment reliability, this means that not all adversarial benchmarks serve as calibration stress tests: ANLI-style model-in-loop benchmarks may give false assurance about calibration robustness, while human-adversarial benchmarks (AdvGLUE) are the higher-risk calibration setting.

The adversarial construction method distinction also provides a conceptual bridge between our NLP findings and the vision-domain calibration literature. \citet{minderer2021revisiting} demonstrated that image corruption (ImageNet-C) consistently degrades calibration --- an observation that motivated our initial hypothesis of universal $\Delta$ECE $> 0$. The key difference is that image corruptions are not model-aware: they do not select examples that the model was already failing on. AdvGLUE human-adversarial examples share this property. ANLI model-in-loop examples do not --- they were selected to cause errors in the model, so the model is already uncertain on them, producing calibration improvement rather than degradation.

\textbf{Finding 2: RLHF alignment exhibits a benchmark-type-conditional calibration effect.} The alignment tax on AdvGLUE (chat worse than base: $\Delta\Delta$ECE $= -0.026$) and the consistent RLHF moderation on ANLI (100\% of rounds, $\Delta\Delta$ECE up to $+0.147$) are not contradictory --- they reflect the same underlying interaction. RLHF-aligned models are trained to produce confident, helpful outputs, which reduces calibration degradation when adversarial examples probe uncertainty (ANLI model-in-loop) but may increase it when adversarial examples exploit decisive wrong answers (AdvGLUE human-adversarial). This finding has practical implications for RLHF research: calibration evaluations conducted on clean benchmarks do not predict adversarial calibration behavior, and benchmark construction method must be specified when reporting RLHF calibration results.

\textbf{Finding 3: Task type is an independent calibration moderator.} AdvGLUE QQP (binary, human-adversarial) shows $\Delta$ECE $= -0.029$ despite having the same construction method as AdvGLUE MNLI ($\Delta$ECE $= +0.071$). This comparison within the same adversarial benchmark isolates task type as an independent variable: 3-class NLI is more susceptible to adversarial calibration degradation than binary classification.

\subsection{Limitations}

\textbf{L1: Single-model evaluation for primary calibration experiments.} Our core calibration results (RQ1--RQ3) use Llama-2-7b-hf only. The 4-model grid originally planned was not completed due to computational constraints. We frame our results as a pilot study establishing the conditional structure within a single model family.

\textit{Why acceptable:} Existence (RQ1) and mechanism validity (RQ2) are model-agnostic findings within the single-model scope. The conditional pattern is internally consistent across all measured cells.

\textbf{L2: Shared clean baseline for ANLI.} All ANLI adversarial cells share one GLUE MNLI clean baseline ($n=200$, ECE $= 0.279$). Llama-2-7b-hf's clean MNLI ECE (0.279) is above the \citet{kadavath2022language} range for clean LLMs (0.05--0.15), reflecting weak NLI accuracy (0.365). Consequently, ANLI R1/R2 $\Delta$ECE values may be partially driven by an inflated clean baseline.

\textit{Why acceptable:} GLUE MNLI is the methodologically correct clean counterpart for ANLI. AdvGLUE MNLI uses the same baseline and shows the expected positive $\Delta$ECE, so baseline inflation does not fully explain the ANLI calibration improvement.

\textbf{L3: BBH-MC commonsense reasoning not tested.} BIG-Bench Hard multiple-choice adversarial variants were excluded for scope reasons. All claims are bounded to NLI and binary classification.

\textbf{L4: RLHF moderation conditionality discovered post-hoc.} The benchmark-type $\times$ RLHF interaction was not predicted in the original H-C1 design --- it emerged when H-C1-V2 reported the AdvGLUE reversal as an unexpected finding. Further experimental validation is needed.

\textbf{L5: Pre-specified quantitative thresholds not met.} Mean confidence on incorrect adversarial predictions was 0.616 --- below the originally hypothesized $\geq$0.70 threshold. Only 1 of 5 adversarial cells exceeded $\Delta$ECE $> 0.05$ (AdvGLUE MNLI); the remaining 4 cells were below this threshold. Adversarial calibration degradation in LLMs is smaller in magnitude than vision-domain analogues suggested.

\subsection{Broader Impact}

This work has positive implications for LLM deployment reliability assessment. Practitioners who test model reliability under adversarial conditions by measuring accuracy robustness are missing the calibration dimension. Adversarial calibration audit ($\Delta$ECE measurement) provides a complementary reliability check that should accompany adversarial accuracy evaluation.

The finding that RLHF alignment can exacerbate calibration degradation on static human-adversarial benchmarks (the alignment tax) is a potential concern for high-stakes deployment of RLHF-aligned models. Systems using chat/instruct variants for NLI or classification tasks may be more susceptible to adversarial calibration degradation than base models, in settings where the adversarial examples were designed to exploit base model behavior.

We do not foresee significant negative impacts of this work beyond the usual concern that measurement methodology can be misused. Our methods are fully described and our code infrastructure is shareable.
